% !TEX root = ../main.tex
% Fonte completa: docs/use-cases/UC-15-creare-consultare-report.md
% Scheda compatta; per i test sono riportate le classi, non tutti i metodi.
\clearpage
\subsection{UC-15 - Creare e consultare report}
\label{uc:UC-15-creare-consultare-report}
\begingroup
\fontsize{9.2}{10.6}\selectfont
\setlength{\tabcolsep}{5pt}
\setlength{\extrarowheight}{2pt}
\arrayrulecolor{lineeUC}
\setlength{\arrayrulewidth}{0.3pt}
\renewcommand{\arraystretch}{1.12}
\setlength{\parskip}{0pt}
\noindent\begin{tabularx}{\textwidth}{>{\columncolor{etichetteUC}\raggedright\arraybackslash\sffamily\bfseries\color{black!75}}p{2.5cm}>{\raggedright\arraybackslash}X}
\hline
ID e livello & UC-15 -- User goal \\[2pt] \hline
Attori & \textbf{Principale:} Administrator; \textbf{secondari:} User segnalato \\[2pt] \hline
Descrizione & Un Administrator registra un report relativo a uno User ordinario, eventualmente collegato a un evento, e consulta i report tramite diversi filtri.\par\textbf{Trigger:} L'Administrator richiede di registrare o consultare un report. \\[2pt] \hline
Precondizioni & \textbf{1.} L'Administrator indicato esiste e possiede privilegi amministrativi. \par \textbf{2.} Lo User segnalato esiste e non è Administrator. \\[2pt] \hline
Postcondizioni & \textbf{1.} In creazione, il report è persistito con severità, commento, riferimenti e data assegnata dal sistema. \par \textbf{2.} In consultazione, il sistema restituisce i report corrispondenti al filtro senza modificare dati. \\[2pt] \hline
Garanzia minima & Un report non valido non viene persistito; una consultazione non produce scritture. \\[2pt] \hline
Flusso principale & \begin{minipage}[t]{\linewidth}\begin{enumerate}[leftmargin=1.4em,nosep,label=\arabic*.]
\item L'Administrator seleziona lo User da segnalare.
\item Indica severità, eventuale commento ed eventuale evento correlato.
\item Il sistema verifica che il chiamante sia Administrator e coincida con l'autore indicato.
\item Il sistema verifica che il bersaglio sia uno User ordinario.
\item Se presente, il sistema verifica l'esistenza dell'evento.
\item Il sistema valida severità e commento.
\item Il sistema assegna la data corrente e persiste il report.
\item Il sistema restituisce l'identificativo generato.
\end{enumerate}\vspace{3pt}\end{minipage} \\[2pt] \hline
Flussi alternativi & \textbf{1a. Consultazione dei report} \par (1) L'Administrator sceglie un criterio: identificativo, User, Administrator autore, evento, severità o loro combinazione supportata. (2) Il sistema verifica che il chiamante sia Administrator ed esegue una transazione read-only. (3) Il sistema restituisce i risultati corrispondenti. \par\smallskip \textbf{3a. Autore inesistente o non Administrator} \par (1) Il sistema segnala l'autore mancante oppure nega l'operazione per ruolo insufficiente. (2) Nessun report viene inserito. \par\smallskip \textbf{4a. Bersaglio inesistente o Administrator} \par (1) Il sistema segnala lo User mancante oppure rifiuta la segnalazione di un Administrator. (2) Nessun report viene inserito. \par\smallskip \textbf{5a. Evento correlato inesistente} \par (1) Il sistema segnala che l'evento indicato non esiste. (2) Il caso d'uso termina senza persistenza. \par\smallskip \textbf{6a. Contenuto non valido} \par (1) Il sistema rileva severità assente oppure commento superiore a 1000 caratteri. (2) La richiesta viene rifiutata. \\[2pt] \hline
Regole &\par\smallskip\textbullet{} \textbf{BR-29:} un report è creato da un Administrator, riguarda uno User ordinario e può riferirsi facoltativamente a un evento esistente.\par\smallskip\textbullet{} La severità è obbligatoria; il commento facoltativo non può superare 1000 caratteri. \\[2pt] \hline
Test & \texttt{Secondary\allowbreak{}Service\allowbreak{}Workflow\allowbreak{}Test}; \texttt{Secondary\allowbreak{}Service\allowbreak{}Validation\allowbreak{}Test}; \texttt{Service\allowbreak{}Query\allowbreak{}Delegation\allowbreak{}Test}.\par Creazione, consultazione, aggiornamento ed eliminazione ricevono separatamente l'identità del chiamante e sono riservate agli Administrator. In creazione, il chiamante deve inoltre coincidere con l'autore registrato nel report. \\[2pt] \hline
\end{tabularx}\par
\endgroup
